% TOP_PROBLEM: {"id": 4800033, "title": "Multiple ergodic averages — Problem 33", "category_id": 11, "category": {"id": 11, "name": "dynamical_systems", "display_name": "Dynamical Systems", "description": "Problems about long-term behavior of deterministic systems, Hamiltonian dynamics, and periodic orbits.", "slug": "dynamical-systems", "order_index": 11, "created_at": "2026-07-31T15:26:25.671Z"}, "status": "open", "problem_number": "AMR-047-0033", "set_id": 13}
% FIRST_POSTED: 2026-10-01
% ENTRY_KIND: statement_only
% UnsolvedMath ID 4800033; source catalogue code AMR-047-0033
% !TeX program = lualatex
% Definition and sources only; this file is not an LLM solution attempt.
\documentclass[11pt,a4paper]{article}
\usepackage[margin=25mm]{geometry}
\usepackage{fontspec}
\setmainfont{lmroman10-regular.otf}[BoldFont=lmroman10-bold.otf,
 ItalicFont=lmroman10-italic.otf,BoldItalicFont=lmroman10-bolditalic.otf]
\usepackage{amsmath,amssymb,mathtools}
\usepackage{unicode-math}
\setmathfont{latinmodern-math.otf}
\AtBeginDocument{\let\setminus\smallsetminus}
\usepackage{xurl,hyperref}
\hypersetup{colorlinks=true,urlcolor=blue}
\setcounter{secnumdepth}{0}
\setlength{\parindent}{0pt}
\setlength{\parskip}{5pt}
\setlength{\emergencystretch}{3em}
\begin{document}
\section*{Multiple ergodic averages — Problem 33}\label{4800033}
{\small UnsolvedMath ID 4800033 · AMR-047-0033 · Dynamical Systems · AMR Open Problem Lists}
\subsection{Definitions and mathematical statement}
Fix distinct exponents $a,b\in(0,1)$. Let
$\xi_1,\eta_1,\xi_2,\eta_2,\ldots$ be mutually
independent Bernoulli variables with
$\mathbb P(\xi_k=1)=k^{-a}$ and
$\mathbb P(\eta_k=1)=k^{-b}$.
Write $u_n(\omega)$ and $v_n(\omega)$ for
the increasing enumerations of the two selected
sets. Both enumerations are infinite almost surely.

The target is one event of probability one
on which the following holds for every
probability space $(X,\mathcal B,\mu)$, commuting
invertible measure-preserving $T,S$, and bounded
functions $f,g$:
\[
 \frac1N\sum_{n=1}^N
       (f\circ T^{u_n(\omega)})(g\circ S^{v_n(\omega)})
       \longrightarrow
       \mathbb E_\mu(f\mid\mathcal I_T)\cdot
       \mathbb E_\mu(g\mid\mathcal I_S).
\]
Here $\mathcal I_T$ is the completed
$\sigma$-algebra of sets invariant modulo $\mu$
under $T$, and similarly for $S$.
Prove convergence in $L^2(\mu)$ and, as
the stronger version, pointwise $\mu$-almost
everywhere.

The same index $n$ pairs the $n$th selected
time of each process; the selection indicators
are not paired by their original locations.
The exponents are fixed in advance, and
the full-probability event may depend on
them but not on the dynamical system.
For pointwise convergence, the dynamical null
set may depend on the chosen observables.
No ergodicity assumption is imposed, so
the asserted limit is a product of conditional
expectations rather than necessarily a constant.
\subsection{Short English statement}
Do independently sampled time sequences with different growth exponents force double ergodic averages to separate into the two invariant means?
\subsection{Sources}
\begin{itemize}
\item[S1] ULAM.AI and the UnsolvedMath Contributors, supplied problems.json, record 4800033 (AMR-047-0033), AMR Open Problem Lists. Catalogue identity and source transcription; CC BY 4.0.
\url{https://huggingface.co/datasets/ulamai/UnsolvedMath}.
\item[S2] Frantzikinakis - Some open problems on multiple ergodic averages (2016), Problem 33. Original source indexed by AMR-047-0033.
\url{https://arxiv.org/abs/1103.3808}.
\end{itemize}
\end{document}
